\documentclass[10pt,a4paper]{article}

\usepackage[utf8x]{inputenc}
\usepackage[T1]{fontenc}
\usepackage[francais]{babel}
\usepackage{lmodern}
\usepackage{amsfonts,mathrsfs}

\usepackage{amsmath}
\usepackage{amssymb}
\usepackage{mathtools}

\usepackage[hmargin=2cm, vmargin=2cm]{geometry}

\usepackage{tabularx}

\usepackage{fancyhdr,lastpage}

\usepackage{hyperref}
\hypersetup{
	pdfauthor={Xavier TISSERAND},     % author
  	}

%%%%%%%%%%%%%%%%%%%%%%%%%%%%%%%%%%%%%%%%%%
%%            Intervalle             %%%%%
%%%%%%%%%%%%%%%%%%%%%%%%%%%%%%%%%%%%%%%%%%
\newcommand{\inter}[4]
{\mathchoice
{\left#1#2\mathclose{}\mathpunct{};#3\right#4}% mode \displaystyle
{\mathopen{#1}#2\mathclose{}\mathpunct{};#3\mathclose{#4}}% mode \textstyle
{\mathopen{#1}#2\mathclose{}\mathpunct{};#3\mathclose{#4}}% mode \scriptstyle
{\mathopen{#1}#2\mathclose{}\mathpunct{};#3\mathclose{#4}}% mode
%\scriptscriptstyle
}

\newcommand{\interff }[2]{\ensuremath{\inter{[}{#1}{#2}{]}}} %ferme ferme
\newcommand{\interof }[2]{\ensuremath{\inter{]}{#1}{#2}{]}}} %ouvert ferme
\newcommand{\interfo }[2]{\ensuremath{\inter{[}{#1}{#2}{[}}} %ferme ouvert
\newcommand{\interoo }[2]{\ensuremath{\inter{]}{#1}{#2}{[}}} %ouvert ouvert



%%%%%%%%%%%%%%%%%%%%%%%%%%%%%%
%%     grands crochets      %%
%%%%%%%%%%%%%%%%%%%%%%%%%%%%%%
\newcommand{\cro}[1]{\left[#1\right]}

%%%%%%%%%%%%%%%%%%%%%%%%%%%%%%
%%     grandes parenthèses  %%
%%%%%%%%%%%%%%%%%%%%%%%%%%%%%%
\newcommand{\pa}[1]{\left(#1\right)}

%e de l'exponentielle
\newcommand{\e}{\mathrm{e}}

%signe inferieur ou egal francais
\renewcommand{\leq}{\leqslant}

%signe superieur ou egal francais
\renewcommand{\geq}{\geqslant}


%%%%%%%%%%%%%%%%%%%%%%%%%%%%%%%%%%%%%%%%%%%%%%%%%%%%%
%Définit un environnement exercice avec un compteur%%
%%%%%%%%%%%%%%%%%%%%%%%%%%%%%%%%%%%%%%%%%%%%%%%%%%%%%
\newcounter{num}
\newcommand{\exo}%
	{%
	\setcounter{Partie}{0}% raz compteur partie
	\addtocounter{num}{1}%
	{%
	\par % commence à une nouvelle ligne
	\noindent\textbf{Exercice~\arabic{num}~:~}%
	}%
	}%

%%%%%%%%%%%%%%%%%%%%%%%%%%%%%%%%%%%%%%%%%%%%%%%%%%%%%
%Définit un environnement partie avec un compteur%%
%%%%%%%%%%%%%%%%%%%%%%%%%%%%%%%%%%%%%%%%%%%%%%%%%%%%%
\newcounter{Partie}
\newcommand{\Partie}[1]%
   {%
   \goodbreak% Evite au tant que possible un titre en bas
	%de page tout seul
   \addtocounter{Partie}{1}
   \noindent\textbf{Partie \Alph{Partie} : #1}%
   \par%
   }% 


\newcommand{\R}{\mathbf{R}}

% lettres grecques à la française
\renewcommand{\epsilon}{\varepsilon} 


%loi normale N
\newcommand{\normale}[2]{\mathscr{N}\pa{#1\mathpunct{};#2}}

%coefficient bonomial
\newcommand{\coeff}[2]{\operatorname{C}_{#1}^{#2}}

%d droit de dx
\newcommand{\dd}{\mathop{}\mathopen{}\mathrm{d}}


%courbe C
\newcommand{\courbec}{\mathscr{C}}

\newcommand{\cad}{c'est-à-dire }

\providecommand{\abs}[1]{\left|#1\right|}


\newcommand{\dl}[2]{développement limité d'ordre {$#1$} en {$#2$}}

\everymath{\displaystyle}


\setlength\parindent{0pt}



%%%%%%%%%%%%%%%%%%%%%%%%%%%%%%%%
%    document tapé sous linux
%%%%%%%%%%%%%%%%%%%%%%%%%%%%%%%%

\title{BTS Groupement B}
\author{Mathématiques}
\date{Session 2012}


\headheight 20.0 pt
\fancyhead[L]{BTS }
\fancyhead[C]{Groupement B}
\fancyhead[R]{2012}
\fancyfoot[C]{{\scriptsize\textsl{ Xavier TISSERAND, Lycée Léonce Vieljeux,
2011-2012}}\\
\scriptsize\thepage/\pageref{LastPage}
}
\pagestyle{fancy}


\begin{document}
\maketitle

\exo \emph{12 points}

\Partie{Résolution d'une équation différentielle}

On considère l'équation différentielle $(E) \ : \quad y'+2y=-5\e^{-2x}$, où $y$
est une fonction inconnue de la variable réelle $x$, définie et dérivable sur
$\R$, et $y'$ la fonction dérivée de $y$.

\begin{enumerate}
 \item Déterminer les solutions de l'équation différentielle $(E_0)\ : y'+2y=0$

\item Soit $g$ la fonction définie sur $\R$ par $g(x)=-5x\e^{-2x}$. Démontrer
que la fonction $g$  est une solution  de $(E)$

\item En déduire les solutions de l'équation différentielle $(E)$.

\item Déterminer la solution $f$ de l'équation différentielle $(E)$ vérifiant
la condition initale $f(0)=1$.
\end{enumerate}

\Partie{Étude locale d'une fonction}

Soit $f$ la fonction définie sur $\R$ par $f(x)=(1-5x)\e^{-2x}$. On note
$\courbec$ sa courbe représentative dans le plan muni d'un repère orthonormal.

\begin{enumerate}
 \item 

\begin{enumerate}
 \item On admet le résultat suivant : $\lim_{x\to +\infty} -5x\e^{-2x}=0$.
Calculer $\lim_{x\to +\infty}f(x)=0$.

\item \emph{Cette question est une question à choix multiples. Une seule
réponse est exacte. Recopier sur la copie la réponse qui vous paraît exacte. On
ne demande aucune justification.}

\emph{La réponse juste rapporte un point. Une réponse fausse ou une absence de
réponse ne rapporte ni n'enlève aucun point.}

La courbe $\courbec$ admet une asymptote en $+\infty$ dont une équation est 

\begin{center}
 \begin{tabularx}{14cm}{|*{3}{>{\centering\arraybackslash$}X<{$}|}}
\hline
y=1-5x&y=0&x=0\\
\hline
 \end{tabularx}
\end{center}
\end{enumerate}

\item 

\begin{enumerate}
 \item \`A l'aide du \dl{2}{0} de la fonction : $t\mapsto \e^t$, déterminer le
\dl{2}{0} de la fonction $x\mapsto \e^{-2x}$.

\item En déduire que le \dl{2}{0} de la fonction $f$ est
\[
f(x)=1-7x+12x^2+x^2 \epsilon (x)\quad \text{avec } \lim_{x\to
0}\epsilon(x)=0
\]

\item En déduire une équation de la tangente $T$ à la courbe $\courbec$ au
point d’abscisse $0$.

\item \emph{Cette question est une question à choix multiples. Une seule
réponse est exacte. Recopier sur la copie la réponse qui vous paraît exacte. On
ne demande aucune justification.}

\emph{La réponse juste rapporte un point. Une réponse fausse ou une absence de
réponse ne rapporte ni n'enlève aucun point.}

On veut justifier qu'au voisinage du point d'abscisse $0$, la courbe $\courbec$
est au-dessus de la droite $T$. Recopier sur la copie la justification qui vous
paraît exacte.

\begin{center}
 \begin{tabularx}{14cm}{|*{3}{>{\centering\arraybackslash}X<{}|}}
\hline
$12x^2$ est positif au voisinage de $0$ & $x^2 \epsilon (x)$ est positif au
voisinage de $0$ & $1-7x$ est positif au voisinage de $0$\\
\hline
 \end{tabularx}
\end{center}
\end{enumerate}

\end{enumerate}

\Partie{Calcul intégral}

\begin{enumerate}
 \item On note $I=\int_1^2 f(x)\dd x$ où $f$ est la fonction définie dans la
partie B.

\begin{enumerate}
 \item Démontrer, à l'aide d'une intégration par parties, que :
$I=\frac{23\e^{-4}-13\e^{-2}}{4}$.

\item Donner une valeur approchée de $I$, arrondie à $10^{-2}$.
\end{enumerate}

\item 

\begin{enumerate}
 \item Donner, sans justification, le signe de $f(x)$ pour $x$ dans
l'intervalle $\interff{1}{2}$.

\item Interpréter graphiquement le nombre $I$.

\emph{Dans cette question, toute trace de recherche même incomplète ou non
aboutie, sera prise en compte.}
\end{enumerate}

\end{enumerate}

\newpage

\exo \emph{8 points}

\begin{center}
 \textbf{Les trois parties de cet exercice peuvent être traitées de façon
indépendante.}
\end{center}

Un particulier souhaite acheter, auprès d'un producteur, des bottes de paille
pour l'isolation de sa maison.

\begin{center}
 \textbf{Dans cet exercice, les résultats approchés sont à arrondir à
}$\mathbf{10^{-2}.}$
\end{center}

\Partie{Loi normale}

On prélève au hasard une botte de paille dans la production du 20 juillet 2011.

\begin{enumerate}
 \item On note $X$ la variable aléatoire qui, à chaque botte ainsi prélevée,
associe son épaisseur, exprimée en millimètres. On admet que $X$ suit la loi
normale de moyenne $300$ et d’écart type $18$.

Calculer la probabilité $p(350\leq X\leq 370)$.

\item On note $Y$, la variable aléatoire qui, à chaque botte prélevée dans la
production de cette journée, associe sa densité exprimée en kg/m$^3$. On admet
que $Y$ suit la loi normale de moyenne $100$ et d’écart type $5$.

Calculer la probabilité qu’une botte, prélevée dans la production de cette
journée ait une densité comprise entre 90 kg/m$^3$ et 110 kg/m$^3$.

\item On suppose que les variables aléatoires $X$ et $Y$ sont indépendantes.

Une botte de paille est conforme aux normes d'isolation si son épaisseur,
exprimée en millimètres, appartient à l'intervalle $\interff{350}{370}$ et si
sa densité, exprimée en kg/m$^3$, appartient à l'intervalle $\interff{90}{110}$.

Calculer la probabilité qu'une botte prélevée dans la production de cette
journée soit conforme aux normes d'isolation.

\emph{Dans cette question, toute trace de recherche même incomplète ou non
aboutie, sera prise en compte.}
\end{enumerate}

\Partie{Loi binomiale}

On considère un stock important de bottes de paille, dont une partie est
destinée à un usage d'isolation. On note $E$ l'événement : \og une botte
prélevée au hasard dans le stock est conforme aux normes d'isolation\fg.

On suppose $p(E)=0,4$.

On prélève au hasard 5 bottes de paille dans le stock pour vérification de la
conformité aux normes d'isolation. Le stock est suffisamment important pour que
l'on puisse assimiler ce prélèvement à un tirage avec remise de $5$ bottes.

On considère la variable aléatoire $Z$ qui, à tout prélèvement de $5$ bottes
ainsi défini, associe le nombre de bottes de paille conformes aux normes
d'isolation.

\begin{enumerate}
 \item Justifier que la variable aléatoire $Z$ suit une loi binomiale dont on
précisera les paramètres.

\item Calculer la probabilité que, dans un tel prélèvement, toutes les bottes
de paille soient conformes aux normes d'isolation.

\item Calculer la probabilité que, dans un tel prélèvement, au moins quatre
bottes de paille soient conformes aux normes d'isolation.
\end{enumerate}

\Partie{Intervalle de confiance}

Dans cette partie, on considère les bottes de paille produites le 22 juillet
2011. On prélève au hasard un échantillon de $50$ bottes de paille dans cette
production. La production est assez importante pour que l'on puisse assimiler
ce prélèvement à un tirage avec remise.

On constate que $37$ bottes de paille de cet échantillon sont conformes aux
normes d'isolation.

\begin{enumerate}
 \item Donner une estimation ponctuelle de la fréquence inconnue $p$ des
bottes de paille de cette production  qui sont conformes aux normes d'isolation.

\item Soit $F$ la variable aléatoire qui, à tout échantillon de $50$ bottes
ainsi prélevé dans cette production, associe la fréquence des bottes de cet
échantillon qui sont conformes aux normes d'isolation.

On suppose que $F$ suit la loi normale de moyenne inconnue $p$ de d'écart type
$\sqrt{\frac{p(1-p)}{50}}$.

\begin{enumerate}
 \item \label{q2ab_ex2} Déterminer un intervalle de confiance de la fréquence
$p$ au niveau de confiance de $95\%$.

\item On considère l’affirmation suivante : \og la fréquence $p$ est
obligatoirement dans l'intervalle de confiance obtenu à la question
\ref{q2ab_ex2}\fg.

Cette affirmation est-elle vraie ? (Donner la réponse sans explication.)
\end{enumerate}

\end{enumerate}


 \vspace{2cm}
Suggestions ou remarques :
\href{mailto:xavier.tisserand@ac-poitiers.fr}{xavier.tisserand@ac-poitiers.fr}

%Xavier TISSERAND est professeur de mathématiques en cpge ATS au lycée Vieljeux
%de La Rochelle
% Cette classe ATS est exclusivement réservée aux étudiant(e)s titulaires
% d'un BTS ou DUT industriel qui veut intégrer une école d'ingénieur, sur
%concours  ou sur dossier. La durée de formation est de un an, sans possibilité
%de  redoublement. Il existe en France 31 classes de ce type. Pour de plus
%amples  renseignements, vous pouvez consulter le site internet du lycée
% 
% http://www.lycee-vieljeux.fr/cpge/ats.html
% 
% et/ou me contacter par mail à l'adresse ci-dessus.


\end{document}
